\section{Data Dictionary and Naming Conventions}
\label{sec:data-dictionary}

\sectionnote{Owner: everyone adds entries; Sydney Eichenberg checks it. Target: 0.5 page. Rubric: Consistency (box names, legends, etc.).}

\begin{requirements}
  \item ``Great'' on this criterion means everything is consistent throughout: box names, legends and terms match in every diagram and in the text.
  \item Both tables start from A1's. Add each new term or convention when you first use it, and update an A1 entry whose meaning changed in Section~\ref{sec:conceptual}.
  \item Final check: every box name in every diagram appears in the Data Dictionary with the same spelling, and every abbreviation used in the report appears in the Naming Conventions.
\end{requirements}

\subsection{Data Dictionary}
\label{sec:dictionary}

\begin{requirements}
  \item Add the A2 analysis terms: concrete architecture, reflexion model, convergence, divergence, absence, crate.
  \item Add any new conceptual components from Section~\ref{sec:conceptual} and the sub-components from Section~\ref{sec:subsystem}.
  \item Add the code elements the report names, with their type (for example \texttt{Mux}, \texttt{Domain}, \texttt{LocalPane}, \texttt{ClientDomain}, \texttt{RemoteSshDomain}).
\end{requirements}

\begin{table}[H]
  \centering
  \caption{Data dictionary.}
  \label{tab:data-dictionary}
  \small
  \begin{tabularx}{\linewidth}{@{}>{\raggedright\arraybackslash}p{0.24\linewidth}>{\raggedright\arraybackslash}p{0.24\linewidth}X@{}}
    \toprule
    \textbf{Term} & \textbf{Type} & \textbf{Definition} \\
    \midrule
    GUI Front End & WezTerm subsystem & Displays terminal content and receives user input through graphical windows. \\
    Multiplexer & WezTerm subsystem & Coordinates windows, tabs, panes, and domains, routing requests between the GUI Front End and the relevant local or remote session. \\
    Mux Client & WezTerm subsystem & Manages connections between the local application and multiplexing domains, sending input and receiving pane updates; for remote SSH domains, it communicates with the remote Mux Server. \\
    Mux Server & WezTerm subsystem & Headless WezTerm instance that manages remote session state, including windows, tabs, panes, and their programs. \\
    Configuration Engine & WezTerm subsystem & Evaluates Lua configuration and supplies runtime settings, key bindings, themes, and event-handling support to the GUI Front End and Multiplexer. \\
    SSH Client & WezTerm subsystem & Establishes the authenticated, encrypted SSH connection used to reach a remote Mux Server. \\
    SSH Server & External system & The remote \texttt{sshd} service that accepts SSH connections and authenticates users. \\
    Terminal Emulation Core & WezTerm subsystem & Interprets terminal output, including escape sequences, and updates a pane's logical screen state. \\
    PTY & Operating-system interface & Pseudo-terminal interface that carries input and output bytes between WezTerm and a shell or other program. \\
    Shell & External program & Command-line program running in a PTY that receives commands and produces terminal output. \\
    Domain & Session/execution context & Determines where a pane's program runs; panes belong to a domain, which may be local or remote. \\
    SSH multiplexing domain & Remote domain type & SSH domain that connects to a remote Mux Server, allowing its panes to keep running independently of the local GUI connection. \\
    Pane & User-interface element & A terminal area within a tab. \\
    Tab & User-interface element & A group of one or more panes within a WezTerm window. \\
    Window & User-interface element & A top-level graphical window that can contain tabs and panes. \\
    predicted echo & Latency feature & A local prediction of typed input, displayed to reduce perceived delay during slower remote sessions. \\
    \bottomrule
  \end{tabularx}
\end{table}

\subsection{Naming Conventions}
\label{sec:naming-conventions}

\begin{requirements}
  \item Replace A1's rule for sequence messages (the last row) with the A2 rule: messages are the concrete method calls, written \texttt{Type::method()}.
  \item Add a row for code identifiers in monospace exactly as in the source: crate names as in \texttt{Cargo.toml} (\texttt{wezterm-client}), file paths (\texttt{mux/src/ssh.rs}) and methods (\texttt{Domain::attach()}).
  \item Add a row for the reflexion diagrams: the line style that marks a convergence, a divergence and an absence.
  \item Add a row saying how a reader matches a conceptual component to its code, pointing to the mapping table in Section~\ref{sec:derivation}.
  \item Add the new component names to the ``Component names'' row.
\end{requirements}

To keep the report and diagrams consistent, we use the following naming conventions:

\begin{table}[!htbp]
  \centering
  \caption{Naming conventions and abbreviations.}
  \label{tab:naming-conventions}
  \begin{tabularx}{\linewidth}{@{}p{0.30\linewidth}X@{}}
    \toprule
    \textbf{Convention} & \textbf{Meaning} \\
    \midrule
    WezTerm & Always capitalize the product name as \textit{WezTerm}. \\
    Component names & Use title case for \textit{GUI Front End}, \textit{Multiplexer}, \textit{Mux Client}, \textit{Mux Server}, \textit{SSH Client}, and \textit{Terminal Emulation Core}. \\
    Mux & Use \textit{Mux} only in the names \textit{Mux Client} and \textit{Mux Server}; call the coordinating subsystem the \textit{Multiplexer}. \\
    Abbreviations & Capitalize \textit{GUI} (Graphical User Interface), \textit{SSH} (Secure Shell), \textit{PTY} (pseudo-terminal), and \textit{OS} (Operating System). \\
    External participants & In sequence diagrams, mark non-WezTerm participants with the UML stereotype \texttt{\guillemotleft{}external\guillemotright{}}, as for the Shell and SSH Server. In prose, identify them as external when relevant. \\
    Generic interface terms & Keep \textit{window}, \textit{tab}, \textit{pane}, \textit{domain}, and \textit{session} lowercase except at the start of a sentence. \\
    Dependency arrows & Solid arrows in the dependency diagram point from a component to a component it depends on; arrow labels state the dependency or flow. \\
    Sequence messages & Use concise, action-oriented lowercase labels, such as \textit{attach domain}, \textit{push changed lines}, and \textit{repaint changed lines}. Preserve capitalization for API or event identifiers such as \texttt{DetachDomain} and \texttt{PaneOutput}. \\
    \bottomrule
  \end{tabularx}
\end{table}
